% answers_short.tex —— 2006 年·北京卷·文科 · 简短答案速查
% 由 tools/make_short_answers.py 自 answers.tex 的 \gksolhead/\gkanswer 提取，
% 并与 2006/2006_数学真题_北京文_答案要点.md 逐题比对后生成，勿手改（改源 MD 后重跑）。

\begin{gknotice}
\textbf{使用说明}：本节只给最终答案便于快速核对，完整推导见解析版。答案为编写者独立推导，\textbf{非官方标准答案}。
\end{gknotice}

\gkpart{选择题}
\begin{center}\begin{tabular}{ccccccccc}
\toprule
题号 & 1 & 2 & 3 & 4 & 5 & 6 & 7 & 8 \\
\midrule
答案 & A & B & C & B & D & B & C & C \\
\bottomrule
\end{tabular}\end{center}

\gkpart{填空题}
\begin{itemize}[leftmargin=2.6em]
  \item[\textbf{9.}] $a=4$。
  \item[\textbf{10.}] 84。
  \item[\textbf{11.}] $a=2$。
  \item[\textbf{12.}] $90^\circ$。
  \item[\textbf{13.}] $a:b:c=5:7:8$，$\angle B=60^\circ$。
  \item[\textbf{14.}] 最小值 $\sqrt2$，最大值 $\sqrt{10}$。
\end{itemize}

\gkpart{解答题}
\begin{itemize}[leftmargin=2.6em]
  \item[\textbf{15.}] (1) $\{x\in\mathbb R:x\ne\frac\pi2+k\pi,\ k\in\mathbb Z\}$；(2) $\dfrac{49}{15}$
  \item[\textbf{16.}] (1) $x_0=1$；(2) $a=2,\ b=-9,\ c=12$
  \item[\textbf{17.}] (1) 证明见解析；(2) $\arccos\dfrac1{\sqrt5}$
  \item[\textbf{18.}] (1) $\dfrac34$；(2) $\dfrac{43}{100}$
  \item[\textbf{19.}] (1) $\dfrac{x^2}{9}+\dfrac{y^2}{4}=1$；(2) $8x-9y+25=0$
  \item[\textbf{20.}] (1) $a_n=22-2n$；(2) $a_n=A+(n-1)d$，其中 $A=a_1,d\in\mathbb Z$，且 $(d\ge0,A\ge6)$，或 $(d=-1,A\ge12)$，或 $(d\le-2,A\ge1-10d)$；共有无穷多个数列。此结论不表示题面有误，也不改动原不等号
\end{itemize}
